\documentclass[../../main.tex]{subfiles}
\begin{document}


{\bf [解]}

% \begin{figure}[htb]
% \centering
% \begin{tikzpicture}[scale=2.6]
%   % 半直線L1,L2 (alpha=25度の例)
%   \def\al{25}
%   \draw[->] (-0.3,0) -- (1.5,0) node[right]{$x$};
%   \draw[->] (0,-0.3) -- (0,1.5) node[left]{$y$};

%   \draw[thick,gray] (0,0) -- ({1.3*cos(\al)},{-1.3*sin(\al)}) node[right]{$L_1$};
%   \draw[thick,gray] (0,0) -- ({1.3*cos(\al)},{1.3*sin(\al)}) node[right]{$L_2$};
%   \coordinate (O) at (0,0);
%   \coordinate (P) at ({0.8*cos(\al)},{-0.8*sin(\al)});
%   \coordinate (Q) at ({0.55*cos(\al)},{0.55*sin(\al)});
%   % R: MQベクトルを-90度回転しsqrt(3)倍してMに加える (高さ=sqrt(3)*MQ=(sqrt(3)/2)PQ となり正三角形を保証)
%   \coordinate (M) at ($(P)!0.5!(Q)$);
%   \coordinate (R) at ($(M)!{sqrt(3)}!-90:(Q)$);
%   \draw[thick] (P) -- (Q) -- (R) -- cycle;
%   \filldraw (O) circle(0.46pt) node[left]{$O$};
%   \filldraw (P) circle(0.46pt) node[below]{$P$};
%   \filldraw (Q) circle(0.46pt) node[above]{$Q$};
%   \filldraw (R) circle(0.46pt) node[right]{$R$};
%   \draw (0.25,0) arc (0:\al:0.25) node[midway,right]{\scriptsize$\alpha$};
%   \draw (0.25,0) arc (0:-\al:0.25);
% \end{tikzpicture}
% \caption{$L_1,L_2$上の$P,Q$と正三角形$PQR$}
% \end{figure}

\paragraph{(1)}

\begin{figure}[htb]
\begin{center}
\begin{tikzpicture}[scale=2.6, >=stealth]

  % --- 1. パラメータ設定 (alpha = 25度) ---
  \def\al{25}
  \pgfmathsetmacro{\tanAl}{tan(\al)}
  % PQ // y軸 かつ |PQ| = 1 より、SQ = SP = 0.5
  % S の x 座標: xS = 0.5 / tan(alpha)
  \pgfmathsetmacro{\xS}{0.5 / \tanAl}
  % R の x 座標: xR = xS + sqrt(3)/2
  \pgfmathsetmacro{\xR}{\xS + sqrt(3)/2}

  % 各頂点・交点座標
  \coordinate (O) at (0, 0);
  \coordinate (S) at (\xS, 0);
  \coordinate (Q) at (\xS, 0.5);
  \coordinate (P) at (\xS, -0.5);
  \coordinate (R) at (\xR, 0);

  % --- 2. 座標軸 ---
  \draw[->, thick] (-0.3, 0) -- ({\xR + 0.5}, 0) node[right] {$x$};
  \draw[->, thick] (0, -0.9) -- (0, 0.9) node[above] {$y$};
  \node[below left, font=\small] at (O) {$O$};

  % --- 3. 半直線 L1, L2 ---
  \draw[thick, gray!80] (O) -- ({\xR + 0.3}, {(\xR + 0.3)*\tanAl}) 
    node[right, black, font=\small] {$L_2: y = x\tan\alpha$};
  \draw[thick, gray!80] (O) -- ({\xR + 0.3}, {-(\xR + 0.3)*\tanAl}) 
    node[right, black, font=\small] {$L_1: y = -x\tan\alpha$};

  % --- 4. 正三角形 PQR の描画（着色ハイライト） ---
  \fill[blue!10] (P) -- (Q) -- (R) -- cycle;
  \draw[ultra thick, blue!80!black] (P) -- (Q) -- (R) -- cycle;

  % 線分 SR (高さ: 点線)
  \draw[thick, dashed, red!80!black] (S) -- (R);

  % --- 5. 角度 alpha の表示 ---
  \draw[->, thin] (0.45, 0) arc (0:\al:0.45);
  \node[font=\scriptsize] at ({\al/2}:0.6) {$\alpha$};
  \draw[->, thin] (0.45, 0) arc (0:-\al:0.45);
  \node[font=\scriptsize] at ({-\al/2}:0.6) {$\alpha$};

  % --- 6. 直角マーク (点 S における PQ と x 軸の交差) ---
  \draw[thin, black] 
    ($(S) + (-0.08, 0)$) -- 
    ($(S) + (-0.08, 0.08)$) -- 
    ($(S) + (0, 0.08)$);

  % --- 7. 寸法線・長さの注記 ---

  % OS = 1 / (2 tan alpha)
  \draw[<->, thin, black] (0, -0.65) -- (\xS, -0.65) 
    node[midway, below, font=\scriptsize] {$OS = \dfrac{1}{2\tan\alpha}$};
  \draw[dashed, gray] (\xS, -0.5) -- (\xS, -0.7);

  % SR = sqrt(3)/2 (正三角形の高さ)
  \draw[<->, thin, red!80!black] (\xS, 0.3) -- (\xR, 0.3) 
    node[midway, above, font=\scriptsize] {高さ $\dfrac{\sqrt{3}}{2}$};
  \draw[dashed, gray] (\xR, 0) -- (\xR, 0.35);

  % --- 8. 各点のプロットとラベル ---
  \fill[black] (P) circle (1.2pt) node[below left, font=\small] {$P$};
  \fill[black] (Q) circle (1.2pt) node[above left, font=\small] {$Q$};
  \fill[black] (S) circle (1.2pt) node[below left, font=\small] {$S$};
  \fill[red!80!black] (R) circle (1.5pt) 
    node[below right, font=\small, text=red!80!black] {$R\left(\dfrac{1}{2\tan\alpha}+\dfrac{\sqrt{3}}{2},\, 0\right)$};
  \draw[dashed] (\xS,0.5) -- (0,0.5) node[left,font=\small] {$\frac{1}{2}$};
  \draw[dashed] (\xS,-0.5) -- (0,-0.5) node[left,font=\small] {$\frac{-1}{2}$};

\end{tikzpicture}
\end{center}
\caption{線分 $PQ$ が $y$ 軸に平行なときの正三角形 $\triangle PQR$ と点 $R$ の位置}
\end{figure}

原点を$O$とする．線分$PQ$と$x$軸の交点を$S$とすると，$\angle SOQ=\alpha$および$|SQ|=\dfrac{1}{2}$だから$S\left(\dfrac{1}{2\tan\alpha},0\right)$となる．$R$は$x$軸上にあって$S$よりも$x$座標が$\dfrac{\sqrt{3}}{2}$大きいから求める座標は
\begin{align}
 R\left(\dfrac{1}{2\tan\alpha}+\dfrac{\sqrt{3}}{2},0\right)
\end{align}
である．

\paragraph{(2)}

与えられた半直線の方程式は
\begin{align}
 \begin{cases}
  L_1:y=-\tan\alpha\cdot x \\
  L_2:y=\tan\alpha\cdot x 
 \end{cases}
\end{align}
である．

$P,Q$は$L_1,L_2$の上にあるから，媒介変数として$|OP|=a>0$, $|OQ|=b>0$を用いて
\begin{align}
\overrightarrow{OP}=a\begin{pmatrix}\cos\alpha\\-\sin\alpha\end{pmatrix},\qquad\overrightarrow{OQ}=b\begin{pmatrix}\cos\alpha\\ \sin\alpha\end{pmatrix}
\end{align}
とおける．題意の条件$|PQ|=1$より，$a,b$に課せられる条件は$\triangle OPQ$に余弦定理を用いて
\begin{align}
 |PQ|^2 &= |OP|^2 + |OQ|^2 -2|OP||OQ|\cos\angle QOP \\
\therefore
 1 &=a^2+b^2-2ab\cos2\alpha \label{eq:1}
\end{align}
である．

\begin{figure}[htb]
\begin{center}
\begin{tikzpicture}[scale=2.6, >=stealth]

  % --- 1. パラメータ設定 (例: a=1.2, b=0.6, alpha=25度) ---
  \def\al{25}
  \def\aVal{1.2}
  \def\bVal{0.6}
  
  % 頂点座標の定義
  \coordinate (O) at (0, 0);
  
  % P = a(cos alpha, -sin alpha), Q = b(cos alpha, sin alpha)
  \coordinate (P) at ({\aVal*cos(\al)}, {-\aVal*sin(\al)});
  \coordinate (Q) at ({\bVal*cos(\al)}, {\bVal*sin(\al)});
  
  % 中点 M = (P+Q)/2
  \coordinate (M) at ($0.5*(P)+0.5*(Q)$);
  
  % 高さベクトル MR = sqrt(3)/2 * ((b+a)sin alpha, -(b-a)cos alpha)
  \pgfmathsetmacro{\sqrtThreeOverTwo}{sqrt(3)/2}
  \coordinate (R) at ($(M) + {\sqrtThreeOverTwo}*({(\bVal+\aVal)*sin(\al)}, {-(\bVal-\aVal)*cos(\al)})$);

  % --- 2. 座標軸 ---
  \draw[->, thick] (-0.3, 0) -- (1.8, 0) node[right] {$x$};
  \draw[->, thick] (0, -0.6) -- (0, 1.2) node[above] {$y$};
  \node[below left, font=\small] at (O) {$O$};

  % --- 3. 半直線 L1, L2 ---
  \draw[thick, gray!80] (O) -- (1.7, {-1.7*tan(\al)}) node[right, black, font=\small] {$L_1$};
  \draw[thick, gray!80] (O) -- (1.7, {1.7*tan(\al)}) node[right, black, font=\small] {$L_2$};

  % --- 4. 正三角形 PQR の描画（着色ハイライト） ---
  \fill[blue!10] (P) -- (Q) -- (R) -- cycle;
  \draw[ultra thick, blue!80!black] (P) -- (Q) -- (R) -- cycle;

  % --- 5. 中点 M 関連の線・ベクトルの描画 ---
  % 線分 OM (破線)
  \draw[thick, dashed, gray] (O) -- (M);
  
  % ベクトル MR (高さ: 点線)
  \draw[thick, dashed, red!80!black, ->] (M) -- (R);

  % --- 6. 角度 alpha の表示 ---
  \draw[->, thin] (0.45, 0) arc (0:\al:0.45);
  \node[font=\scriptsize] at ({\al/2}:0.6) {$\alpha$};
  \draw[->, thin] (0.45, 0) arc (0:-\al:0.45);
  \node[font=\scriptsize] at ({-\al/2}:0.6) {$\alpha$};

  % --- 7. 直角マーク (点 M における PQ と MR の直交) ---
  % PQ 方向の単位ベクトル v1, MR 方向の単位ベクトル v2
  \coordinate (v1) at ($(M)!0.15cm!(Q)$);
  \coordinate (v2) at ($(M)!0.15cm!(R)$);
  \draw[thin, black] 
    (v1) -- ($ (v1) + (v2) - (M) $) -- (v2);
    
  % --- 8. 媒介変数 a, b の長さの表示 (動径 |OP|, |OQ|) ---
  %\draw[<->, thin, blue!80!black] ($(O)!0.3cm!270:(P)$) -- ($(P)!0.3cm!-90:(O)$) 
  %  node[midway, below left, font=\scriptsize] {$|OP|=a$};
  %\draw[<->, thin, blue!80!black] ($(O)!0.3cm!-270:(Q)$) -- ($(Q)!0.3cm!90:(O)$) 
  %  node[midway, above left, font=\scriptsize] {$|OQ|=b$};

  % --- 9. 各点のプロットとラベル ---
  \fill[black] (P) circle (1.2pt) node[below, font=\small] {$P$};
  \fill[black] (Q) circle (1.2pt) node[above, font=\small] {$Q$};
  
  % 中点 M (赤色強調)
  \fill[red!80!black] (M) circle (1.5pt) 
    node[below, xshift=3pt, font=\small, text=red!80!black] {$M$};
    
  % 頂点 R (軌跡上の点)
  \fill[blue!80!black] (R) circle (1.5pt) 
    node[above right, font=\small, text=blue!80!black] {$R(X,Y)$};

  % TODO :: a,bの表示
\end{tikzpicture}
\end{center}
\caption{$L_1,L_2$上の$P,Q$と正三角形$PQR$}
\end{figure}

PQ$の中点を$Mとする．$\overrightarrow{PQ}=\bigl((b-a)\cos\alpha,\,(b+a)\sin\alpha\bigr)$だから，$\overrightarrow{PQ}\perp\overrightarrow{MR}$，$|\overrightarrow{MR}|=\dfrac{\sqrt{3}}{2}|\overrightarrow{PQ}|$および$R$が直線$PQ$に関して$O$と反対側にあることから
\begin{align}
\overrightarrow{MR}=\dfrac{\sqrt{3}}{2}\begin{pmatrix}(b+a)\sin\alpha \\-(b-a)\cos\alpha\end{pmatrix}
\end{align}
となる．

$\overrightarrow{OM}=\bigl(\dfrac{a+b}{2}\cos\alpha,\,\dfrac{b-a}{2}\sin\alpha\bigr)$だから，$\overrightarrow{OR}$は
\begin{align}
 \overrightarrow{OR}
   &=\overrightarrow{OM}+\overrightarrow{MR} \\
   &=\dfrac{1}{2}\begin{pmatrix}\dfrac{a+b}{2}\cos\alpha\\ \dfrac{b-a}{2}\sin\alpha\end{pmatrix}+\dfrac{\sqrt{3}}{2}\begin{pmatrix}(b+a)\sin\alpha\\ -(b-a)\cos\alpha\end{pmatrix} \\
   &=\dfrac{1}{2}\begin{pmatrix}(a+b)(\cos\alpha+\sqrt{3}\sin\alpha)\\H(\sin\alpha-\sqrt{3}\cos\alpha)\end{pmatrix}
\end{align}
である．
よって$R=(X,Y)$とおくと
\begin{align}
 X=\dfrac{a+b}{2}(\cos\alpha+\sqrt{3}\sin\alpha,\qquad Y=\dfrac{b-a}{2}(\sin\alpha-\sqrt{3}\cos\alpha) \label{eq:2}
\end{align}
である．

さてここで
\begin{align}
 \begin{cases}
  \cos\alpha+\sqrt{3}\sin\alpha = 2\sin\bigl(\alpha+\dfrac{\pi}{6}\bigr) \\
  \sin\alpha-\sqrt{3}\cos\alpha = 2\sin\bigl(\alpha-\dfrac{pi}{3}\bigr)
 \end{cases}
\end{align}
は$0<\alpha<\dfrac{\pi}{3}$のときいずれも$0$ではないから\cref{eq:2}より
\begin{align}
 a+b=\dfrac{2X}{\cos\alpha+\sqrt{3}\sin\alpha},\qquad b-a=\dfrac{2Y}{\sin\alpha-\sqrt{3}\cos\alpha} \label{eq:3}
\end{align}
である．
ここで$a^2+b^2=\dfrac{(a+b)^2+(b-a)^2}{2}$，$ab=\dfrac{(a+b)^2-(b-a)^2}{4}$だから，\cref{eq:1}に\cref{eq:3}を代入して
\begin{align}
 1&=\dfrac{(a+b)^2+(b-a)^2}{2}-\dfrac{(a+b)^2-(b-a)^2}{2}\cos2\alpha \\
 2&=(a+b)^2(1-\cos2\alpha)+(b-a)^2(1+\cos2\alpha) \\
 1&=(a+b)^2\sin^2\alpha+(b-a)^2\cos^2\alpha \\
 1&=\left(\dfrac{2X\sin\alpha}{\cos\alpha+\sqrt{3}\sin\alpha}\right)^2 + \left(\dfrac{2Y\cos\alpha}{\sin\alpha-\sqrt{3}\cos\alpha}\right)^2 \\
\therefore
 & \dfrac{X^2}{\left(\dfrac{\cos\alpha+\sqrt{3}\sin\alpha}{2\sin\alpha}\right)^2} + \dfrac{Y^2}{\left(\dfrac{\sin\alpha-\sqrt{3}\cos\alpha}{2\cos\alpha}\right)^2} = 1
\end{align}

これは$X,Y$に関する楕円の方程式だから，$P,Q$が$PQ=1$を保ちながら動くときの点$R(X,Y)$の軌跡はこの楕円の一部である．よって題意は示された．

{\bf [解説]}

(2)の一般の議論から先に初めて，(1)はその特殊な例とした方が解答自体は簡単だろう．
\cref{eq:2}において$a=b$となるケースを考えれば良い．\cref{eq:1}に$a=b$を代入して
\begin{align}
 1 &= 2a^2 -2a^2\cos2\alpha \\
\therefore
 a &= \dfrac{1}{\sqrt{2(1-\cos2\alpha)}} \\
   &= \dfrac{1}{2\sin\alpha}
\end{align}
だから，\cref{eq:2}に代入して
\begin{align}
 \begin{cases}
 X &= a(\cos\alpha+\sqrt{3}\sin\alpha) = \dfrac{\cos\alpha+\sqrt{3}\sin\alpha}{2\sin\alpha} = \dfrac{1}{2\tan\alpha} + \dfrac{\sqrt{3}}{2} \\
 Y &= 0
 \end{cases}
\end{align}
であり，これは解答と一致する．




\end{document}
OB